\documentclass[11pt,a4paper]{article}

\usepackage[margin=25mm]{geometry}
\usepackage{amsmath,amssymb}
\usepackage{booktabs}
\usepackage{hyperref}
\usepackage{xcolor}
\hypersetup{colorlinks=true,linkcolor=blue!50!black,urlcolor=blue!50!black,
            citecolor=blue!50!black}

\newcommand{\doi}[1]{\href{https://doi.org/#1}{doi:#1}}
\newcommand{\localfile}[1]{\\ \hspace*{1em}\texttt{#1}}
\newcommand{\Rgrid}{R}
\newcommand{\Lgrid}{L}
\newcommand{\Spsc}{S_{\mathrm{PSC}}}
\newcommand{\Siface}{S_{\mathrm{IFACE}}}
\newcommand{\Selec}{S_{\mathrm{ELEC}}}

\title{The docking score implemented in this repository}
\author{Differentiable ZDOCK --- \texttt{src/zdock}}
\date{2026-07-25}

\begin{document}
\maketitle

\begin{abstract}
This note states, in closed form, the scoring function evaluated by
\texttt{zdock.score.docking\_score\_elec} and maximised by
\texttt{zdock.search.docking\_search}, and maps every symbol onto the primary
literature: the pairwise shape-complementarity term of Chen \& Weng
(2003)~\cite{ChenWeng2003psc}, the interface atomic-contact-energy pair
potential of Mintseris \emph{et~al.}\ (2007)~\cite{Mintseris2007}, and the
combination with electrostatics of Chen \emph{et~al.}\
(2003)~\cite{Chen2003zdock}, with the earlier grid-based formulation of Chen
\& Weng (2002)~\cite{ChenWeng2002} given only where it explains a symbol that
survives in the code. It also records where the implementation deviates from
the papers, and why.
\end{abstract}

\section{Total score}

Let $R$ denote the receptor (held fixed at the origin) and $L$ the ligand,
placed by a rigid transform $(q,\mathbf{t})$: a rotation $q$ drawn from a
finite quaternion set, followed by a translation $\mathbf{t}$ on the FFT
lattice. All three terms are evaluated by scattering atoms onto a Cartesian
grid of spacing $h$ and contracting grid functions, so that the translational
scan is a single FFT correlation per rotation --- the construction of
Katchalski-Katzir \emph{et~al.}~\cite{Katchalski1992}, used by every ZDOCK
version~\cite{ChenWeng2002,ChenWeng2003psc,Chen2003zdock}. The total is

\begin{equation}
  S(q,\mathbf{t}) \;=\; \alpha\,\Spsc(q,\mathbf{t})
  \;+\; \Siface(q,\mathbf{t})
  \;+\; \beta\,\Selec(q,\mathbf{t}),
  \label{eq:total}
\end{equation}

and \textbf{higher is better}: \texttt{docking\_search} returns the top-$N$
poses by $S$. Section~\ref{sec:signs} explains why keeping all three terms in a
``higher is better'' convention requires an explicit sign flip on $\Siface$.

\paragraph{Grid.}
$h = 1.2$~\AA{} throughout, the value used by every ZDOCK paper
(``A grid spacing of 1.2~\AA{} is used throughout this
study''~\cite{ChenWeng2003psc}). An atom is a \emph{surface} atom if its
solvent-accessible area, computed with a 1.4~\AA{} water probe, exceeds
1~\AA$^2$; otherwise it is a \emph{core} atom (\cite{ChenWeng2003psc},
Methods; same definition in~\cite{ChenWeng2002}).

\section{Shape complementarity: PSC}

\subsection{Grid functions}

Following Chen \& Weng~\cite{ChenWeng2003psc}, Eq.~(3), each molecule is
described by one complex grid function. Write $\mathcal{S}$ and $\mathcal{C}$
for the sets of surface and core atoms, $r_j$ for the van der Waals radius of
atom $j$, and $\mathbf{x}_c$ for the centre of grid cell $c$. Define the
occupancy sets
\begin{align}
  \mathrm{surf}(c) &\iff \exists j\in\mathcal{S}: \lVert \mathbf{x}_c-\mathbf{x}_j\rVert < r_j, \\
  \mathrm{core}(c) &\iff \exists j\in\mathcal{C}: \lVert \mathbf{x}_c-\mathbf{x}_j\rVert < r_j,
\end{align}
with $\mathrm{core}$ taking precedence, and let
$\mathrm{open}(c)\iff \neg\,\mathrm{surf}(c)\wedge\neg\,\mathrm{core}(c)$.
Note that PSC, unlike the earlier GSC, defines \emph{no} solvent-accessible
surface layer: ``any grid point that does not correspond to an atom is in the
open space''~\cite{ChenWeng2003psc}.

The imaginary channel is the same for both partners and encodes the clash
penalty,
\begin{equation}
  \operatorname{Im}[\Rgrid_{\mathrm{PSC}}(c)]
  = \operatorname{Im}[\Lgrid_{\mathrm{PSC}}(c)]
  = \begin{cases}
      \rho   & \mathrm{surf}(c)\\
      \rho^2 & \mathrm{core}(c)\\
      0      & \mathrm{open}(c),
    \end{cases}
  \label{eq:psc-im}
\end{equation}
while the real channel differs between the two and encodes the favourable
atom-pair count,
\begin{align}
  \operatorname{Re}[\Rgrid_{\mathrm{PSC}}(c)] &=
    \begin{cases}
      \bigl|\{\, j \in R : \lVert \mathbf{x}_c-\mathbf{x}_j\rVert < D + r_j \,\}\bigr|
        & \mathrm{open}(c)\\
      0 & \text{otherwise,}
    \end{cases}
  \label{eq:psc-reR}\\[2pt]
  \operatorname{Re}[\Lgrid_{\mathrm{PSC}}(c)] &=
    \bigl|\{\, k \in L : c = \operatorname{nearest}(\mathbf{x}_k) \,\}\bigr|.
  \label{eq:psc-reL}
\end{align}

\subsection{Score}

Eq.~(4) of~\cite{ChenWeng2003psc} takes the \emph{real part} of the
cross-correlation only:
\begin{equation}
  \Spsc(\mathbf{t})
  = \operatorname{Re}\!\Bigl[\textstyle\sum_{c}
      \Rgrid_{\mathrm{PSC}}(c)\,\Lgrid_{\mathrm{PSC}}(c+\mathbf{t})\Bigr]
  = \underbrace{\sum_c \operatorname{Re}[\Rgrid]\operatorname{Re}[\Lgrid]}_{\text{favourable}}
  \;-\; \underbrace{\sum_c \operatorname{Im}[\Rgrid]\operatorname{Im}[\Lgrid]}_{\text{clash penalty}}.
  \label{eq:psc}
\end{equation}
The first sum is exactly the number of receptor--ligand atom pairs within the
cutoff, because $\operatorname{Re}[\Lgrid]$ places one count at each ligand
atom and $\operatorname{Re}[\Rgrid]$ reports how many receptor atoms lie within
$D+r_j$ of that point. The second sum contributes
\begin{center}
\begin{tabular}{lcc}
  \toprule
  overlap & value & $\rho=3.5$ \\
  \midrule
  surface--surface & $-\rho^2$ & $-12.25$ \\
  surface--core    & $-\rho^3$ & $-42.88$ \\
  core--core       & $-\rho^4$ & $-150.06$ \\
  \bottomrule
\end{tabular}
\end{center}
so deeper interpenetration is punished super-linearly, ``with a higher score
indicating better shape complementarity''~\cite{ChenWeng2003psc}.

\paragraph{Why this term dominates.}
On clash-free surface-placed decoys the bare pair count
$\sum_{ij} n_{ij}$ separates near-native from non-native poses with a
Mann--Whitney AUC of $0.998$, i.e.\ it is by far the strongest single signal in
the score. Omitting it --- as this port did until the present revision --- makes
every contact a net penalty and drives the search towards poses that do not
touch at all.

\section{Desolvation: the IFACE pair potential}

Mintseris \emph{et~al.}~\cite{Mintseris2007} replace ZDOCK's averaged ACE
desolvation by a pairwise statistical potential over $k=12$ optimised atom
types. With $n_{ij}$ the number of contacts between receptor atoms of type $i$
and ligand atoms of type $j$, their Eq.~(11) and Eq.~(14) read
\begin{equation}
  e_{ij} = \ln \frac{n_{ij}}{c_{\text{surface-fraction};ij}},
  \qquad
  E_{\mathrm{IFACE}} = \sum_{i=1}^{k}\sum_{j=1}^{k} e_{ij}\, n_{ij}.
  \label{eq:iface}
\end{equation}
On the grid (their Eq.~(13)) the receptor side pre-multiplies the potential,
\begin{equation}
  W_i(c) = \sum_{j} e_{ij}\, H_j(c),
  \qquad
  H_j(c) = \bigl|\{\, m \in R \ \text{of type } j : \lVert \mathbf{x}_c-\mathbf{x}_m\rVert < 6\,\text{\AA} \,\}\bigr|,
\end{equation}
the 6~\AA{} contact sphere being the one that defines $n_{ij}$
in~\cite{Mintseris2007}, Eqs.~(8)--(9),
and the ligand side is the per-type occupancy indicator of their Eq.~(12),
$\Lgrid_i(c)\in\{0,1\}$, so that
\begin{equation}
  \Siface(\mathbf{t})
  = \sigma \sum_{i}\sum_{c} W_i(c)\,\Lgrid_i(c+\mathbf{t}),
  \qquad \sigma = -1 .
  \label{eq:iface-grid}
\end{equation}

\section{Electrostatics}

Chen \& Weng~\cite{ChenWeng2002} write the electrostatic energy as a
correlation between the receptor's electric potential and the ligand's atomic
charges, using CHARMM19 partial charges, with grid points in the receptor core
assigned zero potential "to avoid the contributions from non-physical
receptor-core/ligand contacts"; Chen \emph{et~al.}~\cite{Chen2003zdock},
Eq.~(2), carry the same term into ZDOCK through the imaginary channel. Here it
is kept as a separate real correlation:
\begin{equation}
  V(c) = \sum_{m \in R} \frac{q_m}{\lVert \mathbf{x}_c-\mathbf{x}_m\rVert}
  \quad (\lVert\cdot\rVert < 8\,\text{\AA},\ \mathrm{open}(c)),
  \qquad
  \Selec(\mathbf{t}) = \sum_c V(c)\, Q_L(c+\mathbf{t}).
\end{equation}

\section{Sign conventions}
\label{sec:signs}

The two ingredient papers disagree about which direction is favourable, and
Chen \emph{et~al.}~\cite{Chen2003zdock} say so explicitly:

\begin{quote}\small
``positive PSC scores indicate good shape complementarity \ldots{} whereas ACE
scores can be positive (unfavorable) or negative (favorable) \ldots{} To make
these two scores compatible, we flip the signs of the PSC scores.''
\end{quote}

They resolve it by flipping PSC and \emph{minimising}. Because the FFT search
here selects poses by \texttt{topk}, this implementation flips the other term
instead: $\sigma=-1$ in Eq.~\eqref{eq:iface-grid} turns the ACE-convention pair
potential into a ``higher is better'' quantity, and every term in
Eq.~\eqref{eq:total} then agrees with $\Spsc$. Omitting this flip makes
$\Siface$ anti-correlated with pose quality (measured AUC $0.127$; the negated
term gives $0.873$), and since $\lvert\Siface\rvert \approx \lvert\Spsc\rvert$
it then dominates the total.

\section{Parameters}

\begin{center}
\begin{tabular}{llll}
  \toprule
  symbol & meaning & initial value & source \\
  \midrule
  $h$        & grid spacing              & $1.2$~\AA  & \cite{ChenWeng2003psc}, Methods \\
  $D$        & PSC pair cutoff offset    & $3.6$~\AA  & \cite{ChenWeng2003psc}, ``Optimizing PSC'' \\
  $\rho$     & PSC clash scale           & $3.5$      & \cite{Chen2003zdock}, Eq.~(2) \\
  $\alpha$   & weight on $\Spsc$         & $1.0$      & see below \\
  $\beta$    & weight on $\Selec$        & $3.0$      & \cite{Chen2003zdock}, p.~82 \\
  $e_{ij}$   & $12\times12$ pair potential & IFACE table & \cite{Mintseris2007}, Eq.~(11) \\
  $q_m$      & partial charges (11 types) & CHARMM19 & \cite{ChenWeng2002}, ``Electrostatics'' \\
  $\sigma$   & pair-term sign            & $-1$       & \cite{Chen2003zdock}, p.~81 \\
  \bottomrule
\end{tabular}
\end{center}

$\rho$, $D$, $\alpha$, $\beta$, $e_{ij}$ and $q_m$ are all exposed as
differentiable tensors, so any of them can be optimised; the values above are
the initialisation. Note that $\alpha$ has no counterpart in
Eq.~(2) of~\cite{Chen2003zdock}, where the scale of the PSC term is fixed by
$\rho$ and the term enters with unit weight; the $\alpha = 0.01$ carried in
earlier revisions of this code belongs to the \emph{grid}-based
formulation of~\cite{ChenWeng2002}, Eq.~(6), whose second term is an averaged
ACE energy on a kcal/mol scale rather than a pair count.

\section{Where each equation comes from}

\begin{center}
\begin{tabular}{llll}
  \toprule
  here & quantity & source & source equation \\
  \midrule
  \eqref{eq:total}   & $S = \alpha \Spsc + \Siface + \beta \Selec$
                     & \cite{Chen2003zdock} & Eq.~(2) \\
  \eqref{eq:psc-im}  & clash channel $\rho,\rho^2$
                     & \cite{ChenWeng2003psc} & Eq.~(3), $\operatorname{Im}$ \\
  \eqref{eq:psc-reR} & receptor pair-count channel
                     & \cite{ChenWeng2003psc} & Eq.~(3), $\operatorname{Re}[R]$ \\
  \eqref{eq:psc-reL} & ligand nearest-grid channel
                     & \cite{ChenWeng2003psc} & Eq.~(3), $\operatorname{Re}[L]$ \\
  \eqref{eq:psc}     & $\Spsc = \operatorname{Re}[R\cdot L]$
                     & \cite{ChenWeng2003psc} & Eq.~(4) \\
  \eqref{eq:iface}   & $e_{ij}$, $E_{\mathrm{IFACE}}$
                     & \cite{Mintseris2007} & Eqs.~(11), (14) \\
  \eqref{eq:iface-grid} & grid form of the pair sum
                     & \cite{Mintseris2007} & Eqs.~(12), (13) \\
  $\sigma = -1$      & sign reconciliation & \cite{Chen2003zdock} & p.~81 \\
  $D = 3.6$~\AA      & pair cutoff offset & \cite{ChenWeng2003psc} & ``Optimizing PSC'' \\
  $\rho = 3.5$       & clash scale & \cite{Chen2003zdock} & Eq.~(2) \\
  $\beta = 3$        & electrostatics weight & \cite{Chen2003zdock} & p.~82 \\
  $h = 1.2$~\AA      & grid spacing & \cite{ChenWeng2003psc} & Methods \\
  \bottomrule
\end{tabular}
\end{center}

The four ZDOCK papers are in \texttt{ref/} of this repository (file names in
the bibliography); \cite{Zhang1997ace,Katchalski1992,Basu2016dockq} are cited
for provenance and are not needed to run the code.

\section{Deviations from the papers}

\begin{enumerate}
  \item \textbf{Rotational sampling.} ZDOCK uses evenly distributed Euler angle
    sets --- $1800$ orientations at $\Delta=20^\circ$, $54\,000$ at $6^\circ$,
    $180\,000$ at $4^\circ$ (\cite{ChenWeng2003psc}, ``Rotational Sampling'').
    This implementation draws uniform random quaternions, so for the same count
    the worst-case angular error is larger. Finer is not automatically better:
    \cite{ChenWeng2003psc} report that ``finer angular intervals tend to rank
    the best hits lower than coarser angular intervals'', $\Delta=4^\circ$
    giving the \emph{lowest} success rate for $N_P > 100$.
  \item \textbf{Ligand occupancy in $\Siface$.} Eq.~(12)
    of~\cite{Mintseris2007} is a binary indicator (their Eq.~(14) uses
    ``$\cong$'', acknowledging that the grid form only approximates
    $\sum_{ij}e_{ij}n_{ij}$), which under-counts when two
    ligand atoms of the same type share a cell. That is negligible at
    $h=1.2$~\AA{} (cell volume $1.7$~\AA$^3$) but not at coarser spacings; the
    quantity is $37\%$ smaller at $h=3.0$~\AA.
  \item \textbf{Cell-volume normalisation.} The clash channel counts grid
    cells, so it scales as $h^{-3}$. \texttt{sc\_cell\_volume\_factor}
    rescales it to the $1.2$~\AA{} reference, making the term
    spacing-invariant; it is the identity at the default spacing.
  \item \textbf{DockQ.} Pose quality is scored against~\cite{Basu2016dockq} by
    an atom-level approximation
    without residue or backbone metadata, and without symmetry handling, so
    the $0.23$ acceptance threshold is not identical to canonical CAPRI. Ligand
    RMSD is reported alongside as a cross-check.
\end{enumerate}

\section{Measured effect of completing PSC}

Twelve interface-deleaked PINDER complexes, default parameters (no training),
$1900$ uniform rotations, $h=1.2$~\AA, top-$2000$ search output, nothing
injected into the candidate set. ``Ceiling'' is the DockQ of the best sampled
rotation at the lattice-snapped native translation --- a parameter-free upper
bound on what the search could return; it is $\ge 0.23$ for $100\%$ of these
complexes throughout.

\begin{center}
\begin{tabular}{lrrrrr}
  \toprule
  score implementation & recall & best DockQ & s@1 & s@10 & s@100 \\
  \midrule
  clash penalty only, $h=3.0$~\AA        & $0.0\%$  & $0.107$ & $0\%$  & $0\%$  & $0\%$ \\
  clash penalty only, $h=1.2$~\AA        & $0.0\%$  & $0.066$ & $0\%$  & $0\%$  & $0\%$ \\
  \ \ + pair-term sign flip              & $16.7\%$ & $0.116$ & $0\%$  & $0\%$  & $0\%$ \\
  \textbf{full PSC, Eqs.\ \eqref{eq:psc-reR}--\eqref{eq:psc}}
                                         & $\mathbf{66.7\%}$ & $\mathbf{0.649}$
                                         & $\mathbf{17\%}$ & $\mathbf{42\%}$ & $\mathbf{50\%}$ \\
  \bottomrule
\end{tabular}
\end{center}

\begin{thebibliography}{9}
\bibitem{ChenWeng2003psc}
  R.~Chen and Z.~Weng.
  \emph{A novel shape complementarity scoring function for protein--protein
  docking.}
  Proteins \textbf{51}(3), 397--408 (2003). \doi{10.1002/prot.10334}.
  \localfile{ref/Chen\_Weng\_2003.pdf}
\bibitem{Chen2003zdock}
  R.~Chen, L.~Li and Z.~Weng.
  \emph{ZDOCK: an initial-stage protein-docking algorithm.}
  Proteins \textbf{52}(1), 80--87 (2003). \doi{10.1002/prot.10389}.
  \localfile{ref/Chen\_et\_al\_2003.pdf}
\bibitem{Mintseris2007}
  J.~Mintseris, B.~Pierce, K.~Wiehe, R.~Anderson, R.~Chen and Z.~Weng.
  \emph{Integrating statistical pair potentials into protein complex
  prediction.}
  Proteins \textbf{69}(3), 511--520 (2007). \doi{10.1002/prot.21502}.
  \localfile{ref/Mintseris\_et\_al\_2007.pdf}
\bibitem{ChenWeng2002}
  R.~Chen and Z.~Weng.
  \emph{Docking unbound proteins using shape complementarity, desolvation, and
  electrostatics.}
  Proteins \textbf{47}(3), 281--294 (2002). \doi{10.1002/prot.10092}.
  \localfile{ref/Chen\_Weng\_2002.pdf}
\bibitem{Zhang1997ace}
  C.~Zhang, S.~Vasmatzis, J.~L.~Cornette and C.~DeLisi.
  \emph{Determination of atomic desolvation energies from the structures of
  crystallized proteins.}
  J.\ Mol.\ Biol.\ \textbf{267}(3), 707--726 (1997).
  \doi{10.1006/jmbi.1997.0859}.
\bibitem{Katchalski1992}
  E.~Katchalski-Katzir, I.~Shariv, M.~Eisenstein, A.~A.~Friesem, C.~Aflalo and
  I.~A.~Vakser.
  \emph{Molecular surface recognition: determination of geometric fit between
  proteins and their ligands by correlation techniques.}
  Proc.\ Natl.\ Acad.\ Sci.\ USA \textbf{89}(6), 2195--2199 (1992).
  \doi{10.1073/pnas.89.6.2195}.
\bibitem{Basu2016dockq}
  S.~Basu and B.~Wallner.
  \emph{DockQ: a quality measure for protein--protein docking models.}
  PLoS ONE \textbf{11}(8), e0161879 (2016).
  \doi{10.1371/journal.pone.0161879}.
\end{thebibliography}

\end{document}
